% BHU PYQ -- Professional Exam Template
\documentclass[11pt,a4paper]{article}

\usepackage[a4paper,top=2.5cm,bottom=2.5cm,left=2.8cm,right=2.8cm,headheight=14pt]{geometry}
\usepackage{amsmath,amssymb}
\usepackage[T1]{fontenc}
\usepackage[utf8]{inputenc}
\usepackage{lmodern}
\usepackage{microtype}
\usepackage{fancyhdr}
\usepackage{enumitem}
\usepackage{listings}
\usepackage{hyperref}
\usepackage{lastpage}
\usepackage{parskip}

\hypersetup{colorlinks=true,urlcolor=black,linkcolor=black,
  pdftitle={CS-101: Problem Solving through C -- BHU PYQ}}

\pagestyle{fancy}
\fancyhf{}
\renewcommand{\headrulewidth}{0.4pt}
\renewcommand{\footrulewidth}{0.4pt}
\fancyhead[L]{\small   \textit{Banaras Hindu University}}
\fancyhead[R]{\small   \textit{CS-101/BCS-101: Problem Solving through C}}
\fancyfoot[C]{\small Page  \thepage\ of \pageref{LastPage}}

\lstset{
  basicstyle=\small tfamily,
  frame=single,
  breaklines=true,
  columns=flexible,
  keepspaces=true
}


\newlist{parts}{enumerate}{2}
\setlist[parts,1]{label=(\alph*),leftmargin=*,itemsep=4pt,topsep=3pt}
\setlist[parts,2]{label=(\roman*),leftmargin=*,itemsep=2pt,topsep=2pt}

\newcommand{\pts}[1]{\hfill   \textit{[#1]}}


\begin{document}

\begin{center}
  {\large\bfseries BANARAS HINDU UNIVERSITY}\\[2pt]
  {\normalsize B.Sc. (Hons.) I Semester Examination, 2016-17}\\[6pt]
  \rule{0.8\linewidth}{0.5pt}\\[6pt]
  {\large\bfseries Subject: Computer Science}\\[4pt]
  {\large\bfseries Paper: CS-101 (Problem Solving through C Programming) /}\\[2pt]
  {\large\bfseries BCS-101 (Introduction to Computer Programming Through C)}\\[4pt]
  {\normalsize Time: Three hours \hfill Maximum Marks: 70}
\end{center}

\rule{\linewidth}{0.4pt}

\smallskip



\noindent   \textit{Instructions: Answer any five questions, including Question No. 1, which is compulsory. All undefined symbols have their usual meanings.}

\bigskip

\medskip


\noindent   \textbf{Question 1.} Attempt all questions: \pts{$2 imes9=18$}

\begin{parts}
  \item Define Algorithm. What are its characteristics?
  \item What is the difference between `=' operator and `==' operator?
  \item When is the register storage class most useful? Give a suitable example.
  \item What is the role of a Preprocessor in `C' programming?
  \item Why is the `\&' operator not used with array name in a   \texttt{scanf()} statement?
  \item What do you mean by `$ o$' operator in C? When is it used?
  \item What will be the output of the following code?

\begin{lstlisting}[language=C]
void main()
{
    int a=3, b=4;
    b %= 3+4;
    a *= a+5;
    printf("%d %d", b, a);
}
\end{lstlisting}
  \item Consider the following code segment and write its final output:

\begin{lstlisting}[language=C]
void main()
{
    int x=3, z;
    z = x / ++x;
    printf("x=%d z=%d", x, z);
}
\end{lstlisting}
  \item What is the output of the following code?

\begin{lstlisting}[language=C]
void main()
{
    int a=30, b=0, c;
    c = (a != 10) && (b = 0);
    printf("c = %d", c);
}
\end{lstlisting}
\end{parts}

\medskip


\noindent   \textbf{Question 2.}

\begin{parts}
  \item What is the difference between top-down approach and bottom-up approach of programming? \pts{2}
  \item Draw a flowchart to find the roots of a quadratic equation $ax^2 + bx + c = 0$. \pts{4}
  \item List the rules for naming a variable in C. \pts{2}
  \item Explain the difference between nested if-else clause and else if clause used in C language. \pts{5}
\end{parts}

\medskip


\noindent   \textbf{Question 3.}

\begin{parts}
  \item Define function. What do you mean by function prototype? Is it always compulsory to write function prototype statement while using a function in a `C' program? Justify your answer with one example. \pts{5}
  \item Describe the two ways of passing arguments to function? When do you prefer to use each of them? \pts{4}
  \item Write a program using function to find $m^n$ where $m, n \ge 0$. \pts{4}
\end{parts}

\medskip


\noindent   \textbf{Question 4.}

\begin{parts}
  \item State the difference between automatic and static storage class variables by using appropriate example. \pts{4}
  \item Write your own function for the following operations on string: \pts{6}
  
\begin{parts}
    \item Comparison of two strings
    \item Copy one string into another string
  \end{parts}
  \item What is the use of enumerated data type in `C' programming? \pts{3}
\end{parts}

\medskip


\noindent   \textbf{Question 5.}

\begin{parts}
  \item How are 2D array elements stored in computer memory in `C' programming? \pts{2}
  \item Write a `C' program to find out transpose of a given matrix. \pts{4}
  \item What is meant by the following terms? Give examples of use of each of them. \pts{7}
  
\begin{parts}
    \item Array of pointers
    \item Pointer to Array
  \end{parts}
\end{parts}

\medskip


\noindent   \textbf{Question 6.}

\begin{parts}
  \item Explain the difference between a `Structure' and a `Union' with suitable examples. \pts{4}
  \item How do you pass parameters to main() function in a `C' program? Explain with a suitable example. \pts{4}
  \item Can a structure be used within another structure? Give appropriate example to support your answer. \pts{5}
\end{parts}

\medskip


\noindent   \textbf{Question 7.}

\begin{parts}
  \item What is the need of dynamic memory allocation? Explain the role of   \texttt{realloc()} function with its proper syntax. \pts{5}
  \item Define Macro. In what situation is it better to use Macro instead of using Function? \pts{3}
  \item What is conditional compilation? How does it help a programmer? Explain with an example. \pts{5}
\end{parts}

\vfill
\rule{\linewidth}{0.4pt}

\begin{center}\small
\href{https://bhu-exam.vercel.app/}{Click here to visit our website}\\[4pt]\href{https://me-aryan.vercel.app/}{Click here to visit our contributor}\\[4pt]\href{https://quashan-me.vercel.app/}{Click here to visit our contributor}
\end{center}
\end{document}